\setupcolors[state=start]
\setupcolor[hex]
\definecolor[darkblue][h=003366]
\definecolor[darkgray][h=e6e6e6]
\definecolor[redline][h=d9534f]
\definecolor[blueline][h=5bc0de]
\definecolor[greenline][h=5cb85c]

\setuphead[section][style=\bfa, color=darkblue, before=\blank, after=\blank]
\setuphead[subsection][style=\bf, color=black, before={\blank[2*big]}, after=\blank]
\setuppapersize[A4][A4]
\setupbodyfont[11pt]
\setuplayout[width=middle, height=middle, topspace=2cm, bottomspace=2cm]

\starttext

%%%%%%%%%%%%%%%%%%%%%%%%%%%%%%%%%%%%%%%%%%%%%%%%%%%%%%%%%%%%%%%%%%%%
% ETAPA 1
%%%%%%%%%%%%%%%%%%%%%%%%%%%%%%%%%%%%%%%%%%%%%%%%%%%%%%%%%%%%%%%%%%%%

\title{\bf Modelo OPP - Etapa 1: Ruteo por lote (independiente)}
\blank[big]

\startsection[title={Introducción}]
Este modelo resuelve, para todos los lotes a la vez de forma independiente, el camino de menor tiempo (traslado + atención) que visita las ubicaciones requeridas. Exporta el tiempo de ruta para la Etapa 2.
\stopsection

\startsection[title={Conjuntos}]
\startformula
\startalign
    \NC \mathcal{N} \NC = \text{Órdenes de clientes} \NR
    \NC \mathcal{V} \NC = \text{Locaciones del almacén, incluido } v_{start} \text{ y } v_{end} \NR
    \NC \mathcal{B} \NC = \text{Lotes factibles de pedidos} \NR
    \NC \mathcal{E} \NC = \text{Arcos válidos del grafo de locaciones} \NR
\stopalign
\stopformula
\stopsection

\startsection[title={Parámetros}]
\startformula
\startalign
    \NC v_{start} \NC = \text{Nodo depósito inicial} \NR
    \NC v_{end} \NC = \text{Nodo final} \NR
    \NC T^{travel} \NC = \text{Tiempo por unidad de distancia} \NR
    \NC D_{ij} \NC = \text{Distancia entre locaciones } i,j \in \mathcal{V} \text{ con } (i,j) \in \mathcal{E} \NR
    \NC S_{i} \NC = \text{Tiempo de servicio en la locación } i \in \mathcal{V} \NR
    \NC A_{ni} \NC = \startmathcases
                    \NC 1 \MC \text{si la orden } n \in \mathcal{N} \text{ requiere visitar } i \in \mathcal{V} \NR
                    \NC 0 \MC \text{en otro caso} \NR
                 \stopmathcases \NR
    \NC B_{nb} \NC = \startmathcases
                    \NC 1 \MC \text{si la orden } n \in \mathcal{N} \text{ pertenece al lote } b \in \mathcal{B} \NR
                    \NC 0 \MC \text{en otro caso} \NR
                 \stopmathcases \NR
    \NC L_{b} \NC = \left| \left\{ i \in \mathcal{V} \setminus \{v_{start},v_{end}\} : \sum_{n \in \mathcal{N}} A_{ni} \cdot B_{nb} \ge 1 \right\} \right| \text{ (número de locaciones distintas requeridas)} \NR
\stopalign
\stopformula
\stopsection

\startsection[title={Variables}]
\startformula
\startalign
    \NC z_{ijb} \NC = \startmathcases
                    \NC 1 \MC \text{si el arco } (i,j) \in \mathcal{E} \text{ es usado en el lote } b \in \mathcal{B} \NR
                    \NC 0 \MC \text{en otro caso} \NR
                 \stopmathcases \NR
    \NC q_{ijb} \NC = \text{Flujo de carga (Gavish-Graves): locaciones requeridas del lote } b \in \mathcal{B} \text{ que faltan visitar al recorrer } (i,j) \in \mathcal{E} \NR
    \NC r_{b} \NC = \text{Tiempo de traslado + servicio del lote } b \in \mathcal{B} \text{ bajo la mejor ruta} \NR
\stopalign
\stopformula
\stopsection

\startsection[title={Función objetivo}]
\placeformula
\startformula
\min \sum_{b \in \mathcal{B}} r_{b}
\stopformula
\stopsection

\startsubsection[title={Rutas de lote}]
\startformula
\startalign[n=4, align={right,left,left,left}]
    \NC (1) \NC \sum_{j:(v_{start},j) \in \mathcal{E}} z_{v_{start} j b} = 1; \NC \quad \forall b \in \mathcal{B} \NC \quad \text{(Salida del depósito)} \NR
    \NC (2) \NC \sum_{i:(i,v_{end}) \in \mathcal{E}} z_{i v_{end} b} = 1; \NC \quad \forall b \in \mathcal{B} \NC \quad \text{(Entrada al nodo final)} \NR
    \NC (3) \NC \sum_{i:(i,v_{start}) \in \mathcal{E}} z_{i v_{start} b} = 0; \NC \quad \forall b \in \mathcal{B} \NC \quad \text{(No retorno al depósito)} \NR
    \NC (4) \NC \sum_{j:(v_{end},j) \in \mathcal{E}} z_{v_{end} j b} = 0; \NC \quad \forall b \in \mathcal{B} \NC \quad \text{(No salida desde nodo final)} \NR
\stopalign
\stopformula
\stopsubsection

\startsubsection[title={Visita y conservación de flujo}]
\startformula
\startalign[n=4, align={right,left,left,left}]
    \NC (5) \NC \sum_{i:(i,j) \in \mathcal{E}} z_{i j b} \ge \min\left(1, \sum_{n \in \mathcal{N}} A_{n j} \cdot B_{n b}\right); \NC \quad \forall b \in \mathcal{B},\ j \in \mathcal{V} \setminus \{v_{start},v_{end}\} \NC \quad \text{(Entrada requerida)} \NR
    \NC (6) \NC \sum_{i:(j,i) \in \mathcal{E}} z_{j i b} \ge \min\left(1, \sum_{n \in \mathcal{N}} A_{n j} \cdot B_{n b}\right); \NC \quad \forall b \in \mathcal{B},\ j \in \mathcal{V} \setminus \{v_{start},v_{end}\} \NC \quad \text{(Salida requerida)} \NR
    \NC (7) \NC \sum_{g:(g,i) \in \mathcal{E}} z_{g i b} - \sum_{j:(i,j) \in \mathcal{E}} z_{i j b} = 0; \NC \quad \forall b \in \mathcal{B},\ i \in \mathcal{V} \setminus \{v_{start},v_{end}\} \NC \quad \text{(Conservación en nodos)} \NR
    \NC (8) \NC \sum_{j:(j,i) \in \mathcal{E}} z_{j i b} \le 1; \NC \quad \forall b \in \mathcal{B},\ i \in \mathcal{V} \setminus \{v_{start}\} \NC \quad \text{(A lo sumo un arco de entrada)} \NR
    \NC (9) \NC \sum_{j:(i,j) \in \mathcal{E}} z_{i j b} \le 1; \NC \quad \forall b \in \mathcal{B},\ i \in \mathcal{V} \setminus \{v_{end}\} \NC \quad \text{(A lo sumo un arco de salida)} \NR
\stopalign
\stopformula
\stopsubsection

\startsubsection[title={Subtour elimination y flujo de carga}]
\startformula
\startalign[n=4, align={right,left,left,left}]
    \NC (10) \NC q_{i j b} - L_{b} \cdot z_{i j b} \le 0; \NC \quad \forall b \in \mathcal{B},\ (i,j) \in \mathcal{E} \NC \quad \text{(Flujo solo por arcos usados, Gavish-Graves)} \NR
    \NC (11) \NC \sum_{j:\ (v_{start},j) \in \mathcal{E}} q_{v_{start} j b} = L_{b}; \NC \quad \forall b \in \mathcal{B} \NC \quad \text{(Flujo total desde depósito)} \NR
    \NC (12) \NC \sum_{g:(g,i) \in \mathcal{E}} q_{g i b} - \sum_{j:(i,j) \in \mathcal{E}} q_{i j b} = \min\left(1, \sum_{n \in \mathcal{N}} A_{n i} \cdot B_{n b}\right); \NC \quad \forall b \in \mathcal{B},\ i \in \mathcal{V} \setminus \{v_{start},v_{end}\} \NC \quad \text{(Balance GG)} \NR
\stopalign
\stopformula
\stopsubsection

\startsubsection[title={Tiempo de ruta del lote}]
\startformula
\startalign[n=2, align={right,left}]
    \NC (13) \NC r_{b} = T^{travel} \cdot \left( \sum_{(i,j) \in \mathcal{E}: i \ne v_{end}, j \ne v_{start}} D_{ij} \cdot z_{i j b} \right) \NR
    \NC \NC \quad + \sum_{i \in \mathcal{V} \setminus \{v_{start},v_{end}\}} S_{i} \cdot \left( \sum_{j:(i,j) \in \mathcal{E}} z_{i j b} \right); \quad \forall b \in \mathcal{B} \NR
\stopalign
\stopformula
\stopsubsection


\page
%%%%%%%%%%%%%%%%%%%%%%%%%%%%%%%%%%%%%%%%%%%%%%%%%%%%%%%%%%%%%%%%%%%%
% ETAPA 2
%%%%%%%%%%%%%%%%%%%%%%%%%%%%%%%%%%%%%%%%%%%%%%%%%%%%%%%%%%%%%%%%%%%%

\title{\bf Modelo OPP - Etapa 2: Asignación, secuenciación y cadena de frío}
\blank[big]

\startsection[title={Introducción}]
Recibe el resultado de ruteo de la Etapa 1. Decide la asignación a pickers heterogéneos, secuencia de lotes y controla límite de tiempo fuera de frío restringiendo variables.
\stopsection

\startsection[title={Conjuntos}]
\startformula
\startalign
    \NC \mathcal{N} \NC = \text{Órdenes de clientes} \NR
    \NC \mathcal{P} \NC = \text{Pickers} \NR
    \NC \mathcal{B} \NC = \text{Lotes factibles de pedidos} \NR
\stopalign
\stopformula
\stopsection

\startsection[title={Parámetros}]
\startformula
\startalign
    \NC W^{c}, W^{pen}, W^{e}, W^{t} \NC = \text{Pesos de la función objetivo} \NR
    \NC T^{setup}, T^{pick} \NC = \text{Tiempos de preparación y picking} \NR
    \NC M, M^{seq} \NC = \text{Constantes Big-M} \NR
    \NC C_{n} \NC = \text{Capacidad requerida por la orden } n \in \mathcal{N} \NR
    \NC T^{due}_{n} \NC = \text{Due date de la orden } n \in \mathcal{N} \NR
    \NC B_{nb} \NC = \startmathcases
                    \NC 1 \MC \text{si la orden } n \in \mathcal{N} \text{ pertenece al lote } b \in \mathcal{B} \NR
                    \NC 0 \MC \text{en otro caso} \NR
                 \stopmathcases \NR
    \NC R_{b} \NC = \text{Tiempo de ruta óptimo del lote } b \in \mathcal{B} \text{ (Dato de Etapa 1)} \NR
    \NC V_{p} \NC = \text{Factor de velocidad del picker } p \in \mathcal{P} \NR
    \NC T^{cold} \NC = \text{Tiempo máximo fuera de cámara de frío} \NR
    \NC F_{b} \NC = T^{setup} + T^{pick} \cdot \left( \sum_{n \in \mathcal{N}} C_{n} \cdot B_{nb} \right) \NR
    \NC T^{exec}_{pb} \NC = F_{b} + \frac{R_{b}}{V_{p}} \quad \text{(Tiempo de ejecución como parámetro)} \NR
    \NC \mathcal{PB}^{elig} \NC = \{ (p,b) \in \mathcal{P} \times \mathcal{B} : T^{exec}_{pb} \le T^{cold} \} \NR
\stopalign
\stopformula
\stopsection

\startsection[title={Variables}]
\startformula
\startalign
    \NC x_{pb} \NC = \startmathcases
                    \NC 1 \MC \text{si el lote } b \text{ es asignado al picker } p, \text{ con } (p,b) \in \mathcal{PB}^{elig} \NR
                    \NC 0 \MC \text{en otro caso} \NR
                 \stopmathcases \NR
    \NC c_{pb} \NC = \text{Tiempo de completitud del lote } b \text{ en el picker } p, \text{ con } (p,b) \in \mathcal{PB}^{elig} \NR
    \NC c_{b} \NC = \text{Tiempo de completitud del lote } b \in \mathcal{B} \NR
    \NC y_{bk} \NC = \startmathcases
                    \NC 1 \MC \text{si el lote } b \text{ va antes que } k \text{ en el mismo picker, para } b < k \NR
                    \NC 0 \MC \text{en otro caso} \NR
                 \stopmathcases \NR
    \NC ta_{n}, ea_{n} \NC = \text{Tardanza y adelanto de la orden } n \in \mathcal{N} \NR
\stopalign
\stopformula
\stopsection

\startsection[title={Función objetivo}]
\placeformula
\startformula
\min \left( W^{c} \cdot \sum_{b \in \mathcal{B}} c_{b} + W^{pen} \cdot \left( W^{e} \cdot \sum_{n \in \mathcal{N}} ea_{n} + W^{t} \cdot \sum_{n \in \mathcal{N}} ta_{n} \right) \right)
\stopformula
\stopsection

\startsubsection[title={Asignación y Completitud}]
\startformula
\startalign[n=4, align={right,left,left,left}]
    \NC (1) \NC \sum_{p: (p,b) \in \mathcal{PB}^{elig}} x_{p b} = 1; \NC \quad \forall b \in \mathcal{B} \NC \quad \text{(Asignación única en pares elegibles)} \NR
    \NC (2) \NC T^{exec}_{p b} \cdot x_{p b} \le c_{p b}; \NC \quad \forall (p,b) \in \mathcal{PB}^{elig} \NC \quad \text{(Piso completitud sin Big-M)} \NR
    \NC (3) \NC c_{p b} - M \cdot (1 - x_{p b}) \le c_{b}; \NC \quad \forall (p,b) \in \mathcal{PB}^{elig} \NC \quad \text{(Completitud general del lote)} \NR
\stopalign
\stopformula
\stopsubsection

\startsubsection[title={Secuenciamiento sin solapamiento}]
\startformula
\startalign[n=3, align={right,left,left}]
    \NC (4) \NC c_{p k} \ge c_{p b} + T^{exec}_{p k} - M^{seq} \cdot (3 - x_{p b} - x_{p k} - y_{b k}); \NC \quad \forall p \in \mathcal{P},\ b < k,\ (p,b), (p,k) \in \mathcal{PB}^{elig} \NR
    \NC (5) \NC c_{p b} \ge c_{p k} + T^{exec}_{p b} - M^{seq} \cdot (2 - x_{p b} - x_{p k} + y_{b k}); \NC \quad \forall p \in \mathcal{P},\ b < k,\ (p,b), (p,k) \in \mathcal{PB}^{elig} \NR
\stopalign
\stopformula
\stopsubsection

\startsubsection[title={Earliness y tardiness}]
\startformula
\startalign[n=4, align={right,left,left,left}]
    \NC (6) \NC c_{p b} - T^{due}_{n} - M \cdot (2 - x_{p b} - B_{n b}) \le ta_{n}; \NC \quad \forall n \in \mathcal{N}, (p,b) \in \mathcal{PB}^{elig}: B_{n b} = 1 \NC \quad \text{(Tardanza)} \NR
    \NC (7) \NC T^{due}_{n} - c_{p b} - M \cdot (2 - x_{p b} - B_{n b}) \le ea_{n}; \NC \quad \forall n \in \mathcal{N}, (p,b) \in \mathcal{PB}^{elig}: B_{n b} = 1 \NC \quad \text{(Adelanto)} \NR
\stopalign
\stopformula
\stopsubsection

\startsubsection[title={Dominio de variables}]
\startformula
\startalign[n=3, align={right,left,left}]
    \NC (8) \NC ta_{n}, ea_{n}, c_{b}, c_{p b} \ge 0; \NC \quad \forall n \in \mathcal{N},\ b \in \mathcal{B},\ (p,b) \in \mathcal{PB}^{elig} \NR
    \NC (9) \NC x_{p b}, y_{b k} \in \{0,1\}; \NC \quad \forall (p,b) \in \mathcal{PB}^{elig},\ b < k \NR
\stopalign
\stopformula
\stopsubsection

\stoptext